\documentclass[11pt]{article}
\usepackage[margin=1in]{geometry}
\usepackage{amsmath,amssymb,amsthm,booktabs,microtype}
\usepackage[T1]{fontenc}
\usepackage{lmodern}
\usepackage[hidelinks]{hyperref}
\usepackage{fancyvrb}
\hypersetup{pdftitle={A counterexample to the entropy polynomial root conjecture},pdfsubject={Exact rational counterexample at k=11 and r=10}}
\newtheorem{theorem}{Theorem}
\setlength{\parindent}{0pt}
\setlength{\parskip}{5pt}
\title{A counterexample to the entropy polynomial root conjecture}
\author{}
\date{8 October 2026}
\begin{document}
\maketitle
\vspace{-1.8em}
\begin{abstract}
Wakhare's entropy polynomial root conjecture asserts that an explicit polynomial
$p_{k,r}$ has exactly two roots in $(0,1)$ for every pair of integers $k>r\geq1$.
We give a counterexample with the coprime pair $(k,r)=(11,10)$.
Five exact rational sign evaluations yield at least four distinct roots.
\end{abstract}

\section{The conjecture and the counterexample}
In Conjecture 2 of \cite{Wakhare}, the entropy polynomials are defined by
\begin{equation}\label{eq:h}
h_{k,r}(x)=\sum_{j=0}^{k-1}x^{rj}
 \sum_{v=0}^{j}\frac{(-1)^{j-v}}{v+1}
 \binom{rv+k}{k}\binom{k}{j-v},
\end{equation}
and
\begin{equation}\label{eq:p}
p_{k,r}(x)=\alpha_{k/r}k(1-x^r)^k h_{k,k}(x)
-r(1-x^k)^k h_{k,r}(x),
\end{equation}
where $\alpha_s>0$ is the unique solution of
$\alpha_s(1+\alpha_s)^{s-1}=1$.
The conjecture states that $p_{k,r}$ has exactly two real roots in $(0,1)$,
counting multiplicity.

\begin{theorem}\label{thm:counter}
The polynomial $p_{11,10}$ has a root in each of
\[
(1/5,2/5),\quad (2/5,3/5),\quad (3/5,2/3),\quad (2/3,4/5).
\]
In particular, Conjecture 2 of \cite{Wakhare} is false, even when $k$ and $r$
are required to be coprime.
\end{theorem}

Write $\alpha=\alpha_{11/10}$, so $\alpha^{10}(1+\alpha)=1$.
Define
\begin{align*}
P(t)={}&1+352705t+60632419t^2+1227099358t^3\\
 &+6330005947t^4+10701243741t^5+6330005947t^6\\
 &+1227099358t^7+60632419t^8+352705t^9+t^{10},\\[3pt]
Q(t)={}&11+1939817t+289126442t^2+5380098482t^3\\
 &+25959010187t^4+41238382790t^5+22851341183t^6\\
 &+4098130058t^7+181139618t^8+831413t^9-t^{10}.
\end{align*}
Evaluation of the finite sums in \eqref{eq:h} gives exactly
\begin{equation}\label{eq:coeff}
h_{11,11}(x)=P(x^{11}),\qquad h_{11,10}(x)=Q(x^{10})/11.
\end{equation}
Both $P(t)$ and $Q(t)$ are positive for $0<t<1$: all coefficients of $P$ are
positive, while the constant and final terms of $Q$ sum to $11-t^{10}>0$
and its other coefficients are positive.

\section{Exact signs and the proof}
The function $z^{10}(1+z)$ is strictly increasing on $z>0$. The integer identities
\begin{align*}
117^{10}\cdot242-125^{11}&=-90074807210933370467<0,\\
937^{10}\cdot1937-1000^{11}&=10474898608767871728104442364513>0
\end{align*}
therefore establish
\begin{equation}\label{eq:alpha}
117/125<\alpha<937/1000.
\end{equation}
For $0<x<1$, put
\begin{equation}\label{eq:R}
R(x)=\frac{10(1-x^{11})^{11}Q(x^{10})}
 {121(1-x^{10})^{11}P(x^{11})}.
\end{equation}
The denominator is positive, and \eqref{eq:p}--\eqref{eq:coeff} give
\[
p_{11,10}(x)=11(1-x^{10})^{11}P(x^{11})\bigl(\alpha-R(x)\bigr).
\]
Thus the sign of $p_{11,10}(x)$ is the sign of $\alpha-R(x)$.
The strict rational comparisons below are certified by the positive integers
in the appendix.
\begin{center}
\begin{tabular}{ccc}
\toprule
$x$ & Exact bound on $R(x)$ & $\operatorname{sgn}p_{11,10}(x)$\\
\midrule
$1/5$ & $R(x)<919/1000<117/125$ & $+$\\
$2/5$ & $R(x)>1131/1000>937/1000$ & $-$\\
$3/5$ & $R(x)<935/1000<117/125$ & $+$\\
$2/3$ & $R(x)>938/1000>937/1000$ & $-$\\
$4/5$ & $R(x)<928/1000<117/125$ & $+$\\
\bottomrule
\end{tabular}
\end{center}
Since $p_{11,10}$ is continuous and these five ordered points lie in $(0,1)$,
the intermediate value theorem gives a root in each of the four disjoint
intervals in Theorem~\ref{thm:counter}. The roots are distinct, proving the theorem.
\hfill$\square$

\paragraph{Scope.}
The argument proves \emph{at least four} roots, and makes no assertion that
there are exactly four. It does not refute the associated binary-entropy
inequality, Conjecture 1 of \cite{Wakhare}, proved for all real exponents greater
than one by Ho \cite{Ho}. It also does not resolve Frankl's union-closed sets conjecture.

\paragraph{Reproduction.}
The accompanying \texttt{check\_certificate.py} reconstructs \eqref{eq:coeff}
from \eqref{eq:h}, checks \eqref{eq:alpha}, and verifies all five strict comparisons
using integer arithmetic only. It uses no floating-point arithmetic, computer
algebra system, or numerical root finder. Full numerator and denominator data
are included in \texttt{certificate.json}.

\newpage
\appendix
\section{Cleared denominator certificate}
Write $P(t)=\sum_{j=0}^{10}p_jt^j$ and $Q(t)=\sum_{j=0}^{10}q_jt^j$.
For integers $0<a<b$, define
\begin{align*}
\mathcal P(a,b)&=\sum_{j=0}^{10}p_j a^{11j}b^{11(10-j)},&
\mathcal Q(a,b)&=\sum_{j=0}^{10}q_j a^{10j}b^{10(10-j)},\\
N(a,b)&=10(b^{11}-a^{11})^{11}\mathcal Q(a,b),&
D(a,b)&=121b(b^{10}-a^{10})^{11}\mathcal P(a,b).
\end{align*}
Then $N,D>0$ and $R(a/b)=N(a,b)/D(a,b)$. The following are exact positive
integer values. Each decimal integer is wrapped across lines; its lines are to
be concatenated.

\noindent $919D(1,5)-1000N(1,5)=$
\begin{Verbatim}[fontsize=\footnotesize]
46687482475003686731728494443175716091029147392029835848827061539678
77835881726361080687327974249063360575405536849747865170141728596492
75824805705028730880
\end{Verbatim}

\noindent $1000N(2,5)-1131D(2,5)=$
\begin{Verbatim}[fontsize=\footnotesize]
42538194776507513458711885524826212893895489706184796164361512069476
99426927690709188274361013225559293120108639340014430569971253115788
9439015076159529241505
\end{Verbatim}

\noindent $935D(3,5)-1000N(3,5)=$
\begin{Verbatim}[fontsize=\footnotesize]
52516379104498148655863274985954973433649161277407686847051045991991
85256340399027956110708767358049927938282779151210248609234143827247
884850356820630516531200
\end{Verbatim}

\noindent $1000N(2,3)-938D(2,3)=$
\begin{Verbatim}[fontsize=\footnotesize]
88008107586949258139134527029175262251555372081468008081118660645776
7956959351946751016842608896901672363281250
\end{Verbatim}

\noindent $928D(4,5)-1000N(4,5)=$
\begin{Verbatim}[fontsize=\footnotesize]
12136820388591534024573035987433025741626348214182265217990319570111
14963599358661651488660016203235762863004128345364061271715446024425
596394854034905941742192560
\end{Verbatim}

\begin{thebibliography}{9}
\bibitem{Wakhare}
T.~Wakhare, \emph{Iterated entropy derivatives and binary entropy inequalities},
Journal of Approximation Theory \textbf{307} (2025), 106143.
\href{https://doi.org/10.1016/j.jat.2025.106143}{doi:10.1016/j.jat.2025.106143}.
\href{https://arxiv.org/abs/2312.14743v2}{arXiv:2312.14743v2}.

\bibitem{Ho}
B.~S.~Ho, \emph{A generalization of Boppana's entropy inequality},
\href{https://arxiv.org/abs/2601.19327}{arXiv:2601.19327} (2026).
\end{thebibliography}
\end{document}
